\chapter{Experimental Observations}\label{experimental_observations}
%%% Stokes regime%%%
\section{Stokes regime}

\begin{figure}[!htb]
  \subfloat[\label{subfig-1}]{
    \includegraphics[totalheight=5cm,trim={0 9cm 0 0},clip]{pdf_image_1/ref_15-002.jpg}
  }
  \hfill
  \subfloat[\label{subfig-2:dummy}]{
    \includegraphics[totalheight=4.8cm,trim={1cm 3cm 0 0},clip]{pdf_image_2/ref_16-003.jpg}
  }
     
  \caption{Self-Assembly in stokes flow}
  \label{fig:1}  
  
    
  \end{figure}
  
  Experiments performed by the Migler group\cite{Reference15}\cite{Reference16} show that drop monolayers of an immiscible blend of two polymers(PIB/PDMS) when subjected to shear-flow between two parallel discs give rise to spontaeneously organized "pearl necklace" like structures oriented along the flow direction. In both these studies, the individual components behaved as Newtonian liquids under experimental conditions. PDMS was the dispersed phase and its droplet sizes were of the order of separation between the confining surfaces. Experimental conditions conducive to the formation of these microstructures were explored for a variety of concentrations(1, 5, 10, 20, 28, and 35 $\%$ mass fraction PDMS) and attributed to the composition of the blends, confinement, shear-rate. Cases with even bi-layered drop structures were reported in the limit of certain shear-rates($\dot\gamma$=6-9 $s^{-1}$)\cite{Reference16} and wall separations. It was found that the single-layer necklace-like structures were most stable in lower volume fractions(5-20$\%$). Although the phenomena was hypothesized to be related to the finite-size effects encountered in emulsion rheology: migration of drops from the walls towards the mid-plane and drop collisions, but the lack of a formal understanding of the physics governing the morphology was acknowledged.

\begin{figure}[!htb]
  \subfloat[\label{subfig-1}]{
    \includegraphics[totalheight=3.2cm,trim={0 0 0 0},clip]{pdf_image_3/ref_17-003.jpg}
  }
  \hfill
  \subfloat[\label{subfig-2:dummy}]{
    \includegraphics[totalheight=5cm,angle=90,trim={0 0 0 0},clip]{pdf_image_5/ref_19-003.jpg}
  }
     
  \caption{Formation of string like structures}
  \label{fig:2}  
  
    
  \end{figure}
  
The above mentioned claim that low-density PDMS/PIB blends led to the formation of stable linear microstructures was later supported by studies conducted with PIB/PBMS blends by Vananroye et al.\cite{Reference19}, Peng et al.\cite{Reference17} and with polystyrene(PS) microspheres mixed into a PIB/PDMS blend by Mao et al.\cite{Reference18}.

% In another study Mao et al. studied the role of the asymmetric interfacial affinities of polystyrene(PS) microspheres in PIB/PDMS blends mixed in weight ratios of 10:90 and 90:10 and found that the latter concentration encouraged the encouraged the formation of string-like dropler arrangements in low shear rate(0.1 $s^{-1}$) regimes.

%%% Inertial systems%%%
\section{Inertial regime}

\begin{figure}[!htb]
  
  \subfloat[\label{subfig-1}]{
    \includegraphics[totalheight=3.2cm,trim={0cm 0 0cm 2cm},clip]{pdf_image_1/ref_20-000.jpg}
  }
  %\hfill
  \subfloat[\label{subfig-2:dummy}]{
    \includegraphics[totalheight=5.5cm,angle=90,trim={0 0 0 0},clip]{pdf_image_3/ref_22-007.jpg}
  }
  \\
  \subfloat[\label{subfig-3:dummy}]{
    \includegraphics[totalheight=5cm,trim={0 0 2.82cm 0},clip]{pdf_image_1/ref_20-038.jpg}
  }
  
  \caption{Ordered assembly in hard particles in the presence of inertia}
  \label{fig:3}  
  
    
\end{figure}

Inertial focussing has been of interest particularly due to its ability to localize particles to specific positions in a rectangular channel(including the inertially-driven "tubular pinch" effect\cite{Reference21}) and thus has potential applications in techniques requiring particle or cell separation, filtration, improved encapsulation and even lab-on-chip devices\cite{Reference20}. In a recent study conducted by Humphry et al. it was observed that polystyrene beads($d=9.9 \mu m$) randomly dispersed in a water-surfactant solution(0.1$\%$polysorbate 20) when pressure-driven at a constant speed($0.3m/s$, $\mathit{Re}$=[6.7,10.7])through a microfluidic channel($6cm \times w \times 25\mu m$) gave rise to particle ordering both in axial(long multiparticle trains) and lateral directions\cite{Reference20}. The system was studied for varying aspect-ratios, attained by changing the width of the channel($w=20,40,80,160 \mu m$) and particle volume fractions($\phi = 0.004,0.008,0.016$). As shown in the figure, it was observed that the particles had a tendency to: $a)$ settle in the centre of the longer dimension, $b)$ migrate towards either side of the centerline in the shorter dimension and $c)$ the number of rows(m) in the lateral direction was directly proportional to $w$ and $\phi$. Attempts were made to explain the aforementioned observations. The particle arrangement, with discreet average spacings($l \approx 2d$) attained in a damped oscillatory manner, in the axial direction was attributed to hydrodynamic interactions and the migration along the gradient direction was thought to be due to the presence of open and reversing streamlines in the vicinity of a suspended particle, with additional near-field spiraling regions speculated to be due to inertia. Additionally, a parametric representation involving volume-fraction and channel width was conceived to predict the number of particle-trains in the lateral direction. Exactly similar observations and claims in the case of Poiseuille flow driven by gravity were made by Matas et al.\cite{Reference22}.

%%%Non-Newtonian systems%%%
\section{Non-Newtonian systems}

\begin{figure}[!htb]
  \subfloat[\label{subfig-1}]{
    \includegraphics[totalheight=4.5cm,trim={0 0 2cm 0},clip]{pdf_image_1/ref_23-006.jpg}
  }
  \hfill
  \subfloat[\label{subfig-2}]{
    \includegraphics[totalheight=4.5cm,trim={4.5cm 65.7cm 40cm 25cm},clip]{pdf_image_3/ref_25-003.jpg}
  }
  
  \caption{Self-assembly of RBCs under confinement}
  \label{fig:5}  
  
\end{figure}

Suspensions of living cells are commonly analysed for a variety of purposes like characterization, labelling, sorting and also drug delivery\cite{Reference23}. Many of them are in highly confined geometries having depths of the same order as the particle sizes with the advantage being uniform optical focus that aids in their detection and subsequent analysis and monitoring using microscopy techniques\cite{Reference24}. In an early research Shiga et al. studied a phenomenon of RBC aggregration known as rouleaux formation\cite{Reference25} using a cone-plate viscometer. As seen in the Fig.~\ref{subfig-2}, the aggregration process was discretized into three different stages: a) the first stage, where the one-dimensional aggregates(rouleaux) started to form, b) the second stage, where the rouleaux grew slowly and c) the third stage, where three-dimensional aggregates developed from rouleaux. A few factors affecting the rouleaux formation as mentioned were: decreasing shear rates between $1.9-15 s^{-1}$ to avoid sedimentation on one end and to promote agrregrate formation on the other, RBC hematocrit(volume ratio of $RBC/suspension$) lesser than $0.6 \%$. From a research point of view although the results were presented with clarity, the hypothesis regarding the main attributing factors for rouleaux formation from a particle/cell dynamics perspective was limited.
In a separate study Abkarian et al.\cite{Reference23} studied cellular-level techniques focussing on shape of individual cells in confined geometries and their cross-streamline drift when passed through constrictions. They summarized from past results that trains of RBCs formed in 1-D channels due to the fact that the slower downstream cells slow down the subsequent cells thereby causing intermittent flow of RBSs in smaller capillaries. In an effort to extend the understanding of RBC flow dynamics in pressure-driven flow, they performed experiments in two-dimensional channels with low aspect ratios($height/width$). They attributed the alignment of randomly distributed cells to the a) wall friction which causes their initial deformation due to which they take a parachute-like shape(FIG.~\ref{subfig-1}) and b) confinement which produces long-range inter-particle hydrodynamic interactions.

\begin{figure}[!htb]
  
  \subfloat[\label{subfig-6-1}]{
    \includegraphics[totalheight=5cm,angle=0,trim={0 0 0 0},clip]{pdf_image_2/ref_24-006.jpg}
  }
  \hfill
  \subfloat[\label{subfig-6-2}]{
    \includegraphics[totalheight=4.5cm,angle=0,trim={0 8cm 0 5cm},clip]{pdf_image_2/ref_24-000.jpg}
  }
    
  \caption{System schematics and formation of linear arrays in healthy RBCs and lack of emergence of such order in the case of infected/dead(non-deformable) RBCs}
  \label{fig:6}  
  
\end{figure}


These hypothesis are further reinforced by another study conducted by Iss et al. where they analyzed dynamics for several volume-fractions of RBC suspension under poiseuille flow in two confined geometries(lateral width $w=30 \mu m$ and $60 \mu m$)\cite{Reference24} as shown in Fig.~\ref{fig:6}. They observed that in the vertical direction lift forces localize the deformable RBCs towards a mid-plane thereby creating a 2-D nature of the flow, in the lateral direction too the RBCs move away from the walls and create a cell-free layer due to noncontinuum effects\cite{Reference26} and in the flow direction RBCs self-organized into linear structures. Finally, the linear structures tend to form parallel stacks in along the lateral width especially in the narrower channel($w=30 \mu m$) which had a higher effective volume-fraction . In contrast, rigidified RBCs were observed to not migrate away from the walls and consequently such a system lacked the order as seen in the previous case of deformable RBCs. In conclusion, they note that the spontaeneous formation of stacked linear trains of RBCs is associated with their deformability, vertical and lateral confinements, and the noncontinuum effects of RBCs that cause their additional migration away from wall surfaces(often associated with Fahraeus--Lindqvist effect).